\subsection{Encoding IOR as Temporal Clauses}
\label{sec:case-basics}

The characteristic construction of Section~\ref{sec:trace-observability} applies to every IOR. For structured relations, it reduces to compact laws that look familiar to classical LTL:
\begin{align}
 \text{echo:}\quad
   &G\bigl(y(t)=f(t)\bigr), \label{eq:echo-phi}\\
 \text{one-tick delay:}\quad
   &y(0)=b_0\ \wedge\
     G\bigl(y(t{+}1)=f(t)\bigr), \label{eq:delay-phi}\\
 \text{run length:}\quad
   &y(t)=\phi_{\mathrm{out}}(f,t). \label{eq:rle-phi}
\end{align}
Each equation abbreviates an eligible-output relation $ER(f)$.
Note that equation~\eqref{eq:echo-phi} is memoryless. A Moore system with a fixed initial state cannot make $y(0)$ depend on an input first observed at tick zero, so its functionality cotyledon may be empty. Equation~\eqref{eq:delay-phi} adds the minimum delay required by state transitions and is always realizable by a one-bit register. The fragment therefore exposes the timing constraint already present in Wymore's trajectory semantics.

For Equation~\eqref{eq:rle-phi}, $\phi_{\mathrm{out}}(f,0)=(0,0)$;
tick $2k+1$ emits the bit and run length for symbol $f(2k)$; tick $2k+2$
holds the pair. The next subsection gives the run-length recurrence and the
two-tick pipeline shift. Finite as well as nondeterministic requirements use
the same construction with a restricted time set or a multi-valued $ER(f)$.
